%% kapa: exam=1 round=23 year=2012 period=2 form=A subject=accounting subject_ko=회계학 seq=10 number=10
%% source: Q-Net 감정평가사 기출문제 (cst00309) artlSeq=5052580 file=raw/1008/5052580/2012년 제23회 감정평가사 제1차시험 기출문제(2교시 A형).pdf
%% conversion: kordoc 4.24.51 -> tools/gen1.py (auto); review=proofread-sonnet; flags=none
\begin{question}{10}
회계변경과 오류수정에 관한 설명으로 옳지 않은 것은?
\begin{choices}
  \item 기업은 회계정책의 변경을 반영한 재무제표가 거래, 기타 사건 또는 상황이 재무상태, 재무성과 또는 현금흐름에 미치는 영향에 대하여 신뢰성 있고 더 목적적합한 정보를 제공하는 경우 회계정책을 변경할 수 있다.
  \item 특정 범주별로 서로 다른 회계정책을 적용하도록 규정하거나 허용하는 경우를 제외하고는 유사한 거래, 기타 사건 및 상황에는 동일한 회계정책을 선택하여 일관성 있게 적용한다.
  \item 과거에 발생한 거래와 실질이 다른 거래에 대하여 다른 회계정책을 적용하는 경우는 회계정책의 변경에 해당하지 않는다.
  \item 당기 기초시점에 과거기간 전체에 대한 새로운 회계정책 적용의 누적효과를 실무적으로 결정할 수 없는 경우, 실무적으로 적용할 수 있는 가장 이른 날부터 새로운 회계정책을 소급적용하여 비교정보를 재작성한다.
  \item 전기오류는 특정기간에 미치는 오류의 영향이나 오류의 누적효과를 실무적으로 결정할 수 없는 경우를 제외하고는 소급재작성에 의하여 수정한다.
\end{choices}
\end{question}
\answer{4}
